\begin{center}\begin{tikzpicture}[font=\small,>=stealth,line join=round]
\draw[gray!40,rounded corners=3pt] (0,0) rectangle (14.2,8.3);\node at (7.1,7.700000000000001) {两次偏移抵消，整体也能出现奇偶性};
\begin{scope}[shift={(4.1,3.9499999999999997)},x=1.3043478260869565cm,y=1.0cm]\draw[->,gray] (-2.3,0)--(2.3,0) node[right]{$x$};\draw[->,gray] (0,-2.3)--(0,2.3) node[above]{$y$};\draw[blue,thick,solid,domain=-2:2,samples=180] plot(\x,{\x});\end{scope}\node[anchor=west,align=left,text width=5.0cm] at (8.4,4.0) {$g=x+1$，$f(u)=u-1$\\[10pt]两层本身都非奇非偶\\[10pt]$F=(x+1)-1=x$\\[10pt]整体却为奇函数};
\node[text width=12.6cm,align=center] at (7.1,.65) {口诀给出足够的条件，未穷尽所有复合函数的对称方式。};
\end{tikzpicture}\end{center}
